% This file is part of GUFI, which is part of MarFS, which is released
% under the BSD license.
%
%
% Copyright (c) 2017, Los Alamos National Security (LANS), LLC
% All rights reserved.
%
% Redistribution and use in source and binary forms, with or without modification,
% are permitted provided that the following conditions are met:
%
% 1. Redistributions of source code must retain the above copyright notice, this
% list of conditions and the following disclaimer.
%
% 2. Redistributions in binary form must reproduce the above copyright notice,
% this list of conditions and the following disclaimer in the documentation and/or
% other materials provided with the distribution.
%
% 3. Neither the name of the copyright holder nor the names of its contributors
% may be used to endorse or promote products derived from this software without
% specific prior written permission.
%
% THIS SOFTWARE IS PROVIDED BY THE COPYRIGHT HOLDERS AND CONTRIBUTORS "AS IS" AND
% ANY EXPRESS OR IMPLIED WARRANTIES, INCLUDING, BUT NOT LIMITED TO, THE IMPLIED
% WARRANTIES OF MERCHANTABILITY AND FITNESS FOR A PARTICULAR PURPOSE ARE DISCLAIMED.
% IN NO EVENT SHALL THE COPYRIGHT HOLDER OR CONTRIBUTORS BE LIABLE FOR ANY DIRECT,
% INDIRECT, INCIDENTAL, SPECIAL, EXEMPLARY, OR CONSEQUENTIAL DAMAGES (INCLUDING,
% BUT NOT LIMITED TO, PROCUREMENT OF SUBSTITUTE GOODS OR SERVICES; LOSS OF USE,
% DATA, OR PROFITS; OR BUSINESS INTERRUPTION) HOWEVER CAUSED AND ON ANY THEORY OF
% LIABILITY, WHETHER IN CONTRACT, STRICT LIABILITY, OR TORT (INCLUDING NEGLIGENCE
% OR OTHERWISE) ARISING IN ANY WAY OUT OF THE USE OF THIS SOFTWARE, EVEN IF
% ADVISED OF THE POSSIBILITY OF SUCH DAMAGE.
%
%
% From Los Alamos National Security, LLC:
% LA-CC-15-039
%
% Copyright (c) 2017, Los Alamos National Security, LLC All rights reserved.
% Copyright 2017. Los Alamos National Security, LLC. This software was produced
% under U.S. Government contract DE-AC52-06NA25396 for Los Alamos National
% Laboratory (LANL), which is operated by Los Alamos National Security, LLC for
% the U.S. Department of Energy. The U.S. Government has rights to use,
% reproduce, and distribute this software.  NEITHER THE GOVERNMENT NOR LOS
% ALAMOS NATIONAL SECURITY, LLC MAKES ANY WARRANTY, EXPRESS OR IMPLIED, OR
% ASSUMES ANY LIABILITY FOR THE USE OF THIS SOFTWARE.  If software is
% modified to produce derivative works, such modified software should be
% clearly marked, so as not to confuse it with the version available from
% LANL.
%
% THIS SOFTWARE IS PROVIDED BY LOS ALAMOS NATIONAL SECURITY, LLC AND CONTRIBUTORS
% "AS IS" AND ANY EXPRESS OR IMPLIED WARRANTIES, INCLUDING, BUT NOT LIMITED TO,
% THE IMPLIED WARRANTIES OF MERCHANTABILITY AND FITNESS FOR A PARTICULAR PURPOSE
% ARE DISCLAIMED. IN NO EVENT SHALL LOS ALAMOS NATIONAL SECURITY, LLC OR
% CONTRIBUTORS BE LIABLE FOR ANY DIRECT, INDIRECT, INCIDENTAL, SPECIAL,
% EXEMPLARY, OR CONSEQUENTIAL DAMAGES (INCLUDING, BUT NOT LIMITED TO, PROCUREMENT
% OF SUBSTITUTE GOODS OR SERVICES; LOSS OF USE, DATA, OR PROFITS; OR BUSINESS
% INTERRUPTION) HOWEVER CAUSED AND ON ANY THEORY OF LIABILITY, WHETHER IN
% CONTRACT, STRICT LIABILITY, OR TORT (INCLUDING NEGLIGENCE OR OTHERWISE) ARISING
% IN ANY WAY OUT OF THE USE OF THIS SOFTWARE, EVEN IF ADVISED OF THE POSSIBILITY
% OF SUCH DAMAGE.



\section{Quickstart for macOS}
This guide is a companion to the primary quickstart guide with
information specific to building and using GUFI on macOS. Please refer
to the primary quickstart guide and non-quickstart guides for more
details.

\section{Build From Source}
\subsection{Dependencies}
While macOS comes with some software for building code, much of it is
set up in ways that make building GUFI difficult, if not
impossible. Instead, \texttt{brew} packages should be used. The script
\texttt{contrib/CI/macos.h} should be used to install the
dependencies. The very last line of output from the script contains
extra paths to add to your \texttt{PATH} variable.

\subsection{Configuration}
First create a build directory under the repository root
directory. This guide will create \texttt{build}. \texttt{cd} into it.
\\\\
\texttt{cmake}, which was installed by \texttt{contrib/CI/macos.h},
should be run first to configure the build. \texttt{cmake} tends to
pick up \texttt{AppleClang} for the C and C++ compilers, so explicitly
override them with the environment variables \texttt{CC} and
\texttt{CXX}.

\subsubsection{Additional \texttt{cmake} Flags}
Configuring GUFI on macOS requires one extra flag that is not required
in other environments: \texttt{-DCMAKE\_OSX\_SYSROOT=macosx}. If
\texttt{cmake} is unable to find OpenMP, try also passing in
\texttt{-DOpenMP\_Root="\$(brew --prefix libomp)"}.
\\\\
There are some additional flags that may be used when building on
macOS:
\begin{longtable}{|l|p{0.5\linewidth}|}
  \hline
  \texttt{cmake} Flag & Description \\
  \hline
  \texttt{-DDEP\_AI=<On|Off>} & Enables AI capabilities. Default is
  \texttt{On}, but recommend explicitly turning it on as a reminder
  that it is enabled. \\
  \hline
  \texttt{-DDEP\_INSTALL\_PREFIX=<directory>} & Location to place
  installs of bundled dependencies. Default is \texttt{deps} under
  the \texttt{build} directory, but everything will need to be
  rebuilt if the \texttt{build} directory is wiped. Recommend
  \texttt{\$\{HOME\}/.local}. \\
  \hline
  \texttt{-DDEP\_BUILD\_THREADS=<\#>} & Number of threads to use to
  build dependencies. Default is 1. Recommend setting to higher value.\\
  \hline
\end{longtable}

There are other flags that can be used to change target
directories. Do not change them on macOS.

\subsubsection{Example Configuration}
The basic configuration recommended for macOS is:
\begin{verbatim}
CC="$(brew --prefix llvm)/bin/clang" \
CXX="$(brew --prefix llvm)/bin/clang++" \
cmake .. \
    -DDEP_AI=On \
    -DDEP_INSTALL_PREFIX="${HOME}/.local" \
    -DDEP_BUILD_THREADS="$(nproc --ignore 1)" \
    -DCMAKE_OSX_SYSROOT=macosx
\end{verbatim}

\subsection{Build}
After \texttt{cmake} has succesfully run, simply run \texttt{make}.

\subsection{Convenience Script}
The steps above are also encapsulated in the script
\texttt{contrib/build\_on\_mac.sh}. Run this script from the
repository root. The script was not set up to be sourced, so do not
source it as a shortcut to getting the environment variables.

\subsection{Install}
\subsubsection{Run from \texttt{build}}
GUFI tools may be run from the \texttt{build} directory. Some values
should be added to environment variables:

\begin{itemize}
  \item Add \texttt{\$(realpath
    <DEP\_INSTALL\_PREFIX>/sqlite-vec/lib)} to
    \texttt{DYLD\_LIBRARY\_PATH}
  \item Add \texttt{\$(realpath
    <DEP\_INSTALL\_PREFIX>/sqlite-lembed/lib)} to
    \texttt{DYLD\_LIBRARY\_PATH}
  \item Add \texttt{\$(realpath
    <DEP\_INSTALL\_PREFIX>/llama.cpp/lib)} to
    \texttt{DYLD\_LIBRARY\_PATH}
\end{itemize}

\noindent Optionally,

\begin{itemize}
  \item Add \texttt{\$(realpath build/src)} to \texttt{PATH}
  \item Add \texttt{\$(realpath build/scripts)} to \texttt{PATH}
  \item Add \texttt{\$(realpath build/scripts)} to \texttt{PYTHONPATH}
\end{itemize}

\subsubsection{\texttt{sudo make install}}
\texttt{sudo make install} will install the required files into
\texttt{/usr/local}.
\begin{itemize}
  \item Executables will be placed in
    \texttt{/usr/local/bin}.
  \item Libraries (\texttt{*.dylib}) will be placed
    into \texttt{/usr/local/lib}, including the libraries of the bundled
    dependencies.
  \item If the user facing scripts will be used, the
    configuration file will need to be placed at
    \texttt{/usr/local/etc/GUFI/config}. The example configuration file
    will be placed in \texttt{/usr/local/etc/GUFI/server.example}.
\end{itemize}

\section{Pre-Built Packages}
\subsection{PKGs}
Two pre-built PKGs are attached to the assets of
\href{https://github.com/mar-file-system/GUFI/releases}{each release
 on GitHub}. If proceeding with the PKGs, only download and install
the gufi-server PKG.
\\\\
If installing the PKG via the GUI, double clicking it might result in
a pop-up saying that the action was blocked. If this happens, run
\begin{verbatim}
sudo xattr -r -d com.apple.quarantine <pkg path>
\end{verbatim}
Then double click on it again. Click through the menus. At
``Destination Select", make sure the ``Install for all users of this
computer" option is selected.
\\\\
To install via terminal, run
\begin{verbatim}
sudo installer -pkg <pkg path> -target /
\end{verbatim}
Once the PKG is installed, add \texttt{/usr/local/lib} to
\texttt{DYLD\_LIBRARY\_PATH}.
\\\\
If \texttt{sqlite} is not installed, install it via
\texttt{brew}. Then add \texttt{\$(brew --prefix sqlite)/lib} to
\texttt{DYLD\_LIBRARY\_PATH}. This is necessary because the PKG does
not provide SQLite 3 in order to not potentially overwrite existing
installations and PKGs do not download dependencies themselves.
\\\\
If using the user facing scripts, also add \texttt{/usr/local/lib} to
\texttt{PYTHONPATH}.
\\\\
Note that the PKGs (and other artifacts) on GitHub are specific to
those releases/commits and do not get updated. To get the latest
changes into PKG(s), configure and build GUFI from source. However,
instead of running \texttt{sudo make install}, run
\begin{verbatim}
sudo cpack -G productbuild -DCPACK\_COMPONENTS\_ALL=Server -B <build>
\end{verbatim}
The PKG will be placed in the \texttt{build} directory under the name
\\\texttt{GUFI-<VERSION>-Darwin.pkg}. If \texttt{productbuild} is not
found, run
\begin{verbatim}
xcode-select --install
\end{verbatim}
